\chapter{Brownian Motion}\label{ch:qm:brownian-motion}

A trader buys a stock with an annual volatility of 30\% and places a stop-loss 2\% below the
entry price, then asks the risk desk how likely the stop is to be touched within a month. The
chance that the stock merely \emph{ends} the month below the stop is 41\%; the chance that it
touches the stop at some moment of the month is exactly twice that, 82\%. If the stop is
checked only at each day's close, the answer is 72\%, and a simulation that watches the
closes and forgets what happens between them reports 72\% for a contract that is really
triggered 82\% of the time. The three numbers come from one object, the \omterm{def:qm:brownian-motion:bm}{Brownian motion}, and
from three of its properties: its law at a fixed time, the reflection principle for its
running minimum, and the \omterm{def:qm:brownian-motion:bridge}{Brownian bridge} that describes it between two observations. This
chapter builds the object, proves the properties the series uses, and ends with the
\omterm{def:qm:brownian-motion:qv}{quadratic variation}, the property that makes stochastic calculus different from the
ordinary kind.

\section{From random walks to Brownian motion}

\begin{definition}[Brownian motion, Gaussian process]\label{def:qm:brownian-motion:bm}
A \emph{Gaussian process}\index{Gaussian process} is a process $(X_t)$ whose finite-dimensional
laws $(X_{t_1}, \dots, X_{t_k})$ are all multivariate normal; its law is fixed by its mean
$m(t)$ and covariance $c(s,t)$. A (standard) \emph{Brownian motion}\index{Brownian motion}
$(W_t)_{t \ge 0}$ is an \omterm{def:qm:probability-at-speed:filtration}{adapted process} with $W_0 = 0$, continuous paths, and increments
$W_t - W_s$ independent of $\mathcal F_s$ and distributed $\mathcal N(0, t - s)$ for $s < t$.
\end{definition}

Equivalently, $W$ is the centred \omterm{def:qm:brownian-motion:bm}{Gaussian process} with continuous paths and covariance
$c(s,t) = \min(s,t)$: for $s < t$, $\Cov(W_s, W_t) = \Cov(W_s, W_s + (W_t - W_s)) = s$.
Existence is not obvious, since the definition asks for continuity of uncountably many
random variables at once. Wiener gave the first construction in 1923; the one that serves
simulation builds the path coarse to fine (\cref{sec:qm:brownian-motion:bridge}). The reason
\omterm{def:qm:brownian-motion:bm}{Brownian motion} appears everywhere is a central limit theorem for paths.

\begin{theorem}[Donsker's invariance principle]\label{thm:qm:brownian-motion:donsker}
Let $\xi_1, \xi_2, \dots$ be independent and identically distributed with mean zero and
variance one, $S_k = \sum_{i \le k}\xi_i$, and $X^{(n)}_t = n^{-1/2}(S_{\lfloor nt\rfloor} +
(nt - \lfloor nt\rfloor)\xi_{\lfloor nt\rfloor + 1})$ the linearly interpolated, rescaled walk.
Then $X^{(n)} \xrightarrow{d} W$ as random elements of $C[0, 1]$: $\E[F(X^{(n)})] \to
\E[F(W)]$ for every bounded continuous functional $F$ of the path.
\end{theorem}

\admitted

Donsker's theorem is what licenses the use of \omterm{def:qm:brownian-motion:bm}{Brownian motion} for prices built from many
small, independent trades, whatever their individual law: the maximum of the path, the time
spent below a level, the average price, all continuous functionals, converge together. It
does not license it for prices driven by a few large jumps (chapter~6).

\begin{proposition}[Invariances]\label{prop:qm:brownian-motion:scaling}
If $W$ is a \omterm{def:qm:brownian-motion:bm}{Brownian motion}, so are $-W_t$; $c^{-1/2}W_{ct}$ for every $c > 0$
(\emph{scaling}); $W_{s+t} - W_s$ for fixed $s$; and $tW_{1/t}$ (with value 0 at $t = 0$).
\end{proposition}

\begin{proof}
Each is a centred \omterm{def:qm:brownian-motion:bm}{Gaussian process} with covariance $\min(s,t)$; for scaling, $c^{-1}\min(cs,
ct) = \min(s,t)$; for inversion, $st\min(1/s,1/t) = \min(s,t)$. Continuity at 0 of $tW_{1/t}$
follows from the strong law $W_u/u \to 0$.
\end{proof}

Scaling is the square-root-of-time rule a trader uses without thinking: a stock with 30\%
annual volatility has a one-day standard deviation of $0.30/\sqrt{252} = 1.9\%$ and a
one-month one of $0.30\sqrt{21/252} = 8.7\%$. \Cref{fig:qm:brownian-motion:paths} shows six
paths inside the envelope $\pm 2\sqrt t$.

\begin{omfigure}
\begin{tikzpicture}
\begin{axis}[width=11cm,height=5cm,
  xlabel={$t$ (years)},ylabel={$W_t$},
  tick label style={font=\small},label style={font=\small},
  xmin=0,xmax=1,ymin=-2.4,ymax=2.4,grid=major,grid style={gray!25},
  table/col sep=comma]
\addplot[gray,dashed,thick] table[x=t,y=up]{figdata/methods/02-brownian-motion/paths.csv};
\addplot[gray,dashed,thick] table[x=t,y=dn]{figdata/methods/02-brownian-motion/paths.csv};
\addplot[omDef,thin] table[x=t,y=w0]{figdata/methods/02-brownian-motion/paths.csv};
\addplot[omThm,thin] table[x=t,y=w1]{figdata/methods/02-brownian-motion/paths.csv};
\addplot[omProp,thin] table[x=t,y=w2]{figdata/methods/02-brownian-motion/paths.csv};
\addplot[omExo,thin] table[x=t,y=w3]{figdata/methods/02-brownian-motion/paths.csv};
\addplot[black!60,thin] table[x=t,y=w4]{figdata/methods/02-brownian-motion/paths.csv};
\addplot[omDef!50,thin] table[x=t,y=w5]{figdata/methods/02-brownian-motion/paths.csv};
\end{axis}
\end{tikzpicture}

\omcaption{Six Brownian paths on $[0,1]$ (250 steps) inside the envelope $\pm 2\sqrt t$ (dashed), which
holds 95\% of the law at each fixed time. The spread grows like $\sqrt t$, not like $t$.
Data: the chapter's tutorial, seeded.}\label{fig:qm:brownian-motion:paths}
\end{omfigure}

\section{Path properties and the Markov property}

Brownian paths are continuous and nowhere smooth. Almost surely they are Hölder continuous
of every order $\alpha < \tfrac12$ and of no order $\alpha > \tfrac12$, and differentiable at
no point; they cross any level they touch infinitely often immediately afterwards. None of
this is pathology for its own sake: it is the reason a hedge rebalanced continuously in time
still accumulates a nonzero cost (chapter~3), and the reason a barrier touched once is
touched many times.

\begin{definition}[Markov process, strong Markov property]\label{def:qm:brownian-motion:markov}
An \omterm{def:qm:probability-at-speed:filtration}{adapted process} $(X_t)$ is a \emph{Markov process}\index{Markov process} if for all $s, t
\ge 0$ and bounded measurable $f$, $\E[f(X_{s+t}) \mid \mathcal F_s] = \E[f(X_{s+t}) \mid X_s]$:
given the present, the past adds nothing. It has the \emph{strong Markov property}\index{strong
Markov property} if the same holds with $s$ replaced by any finite \omterm{def:qm:probability-at-speed:stopping}{stopping time} $\tau$.
\end{definition}

\begin{theorem}[Brownian motion is strong Markov]\label{thm:qm:brownian-motion:strongmarkov}
For every \omterm{def:qm:probability-at-speed:stopping}{stopping time} $\tau$ with $\P(\tau < \infty) = 1$, the process $B_t = W_{\tau + t} -
W_\tau$ is a \omterm{def:qm:brownian-motion:bm}{Brownian motion} independent of $\mathcal F_\tau$.
\end{theorem}

\admitted

The proof approximates $\tau$ from above by \omterm{def:qm:probability-at-speed:stopping}{stopping times} with countably many values, for
which the statement is the simple Markov property applied on each value (Karatzas and
Shreve, 1991, §2.6). The \omterm{def:qm:brownian-motion:markov}{strong Markov property} is what makes the next section work: after
the first time a path reaches a level, what it does next is a fresh \omterm{def:qm:brownian-motion:bm}{Brownian motion}, as
likely to go up as down.

\section{The reflection principle and first-passage times}

\begin{definition}[First-passage time]\label{def:qm:brownian-motion:passage}
For a level $b$, the \emph{first-passage time}\index{first-passage time} (or \emph{hitting
time}\index{hitting time}) of a continuous process $X$ is $\tau_b = \inf\{t \ge 0 : X_t = b\}$;
it is a \omterm{def:qm:probability-at-speed:stopping}{stopping time}.
\end{definition}

\begin{theorem}[Reflection principle]\label{thm:qm:brownian-motion:reflection}
For $b > 0$ and $t > 0$, let $M_t = \max_{s \le t}W_s$. Then
\[
\P(M_t \ge b) = \P(\tau_b \le t) = 2\P(W_t \ge b) = 2\bigl(1 - \Phi(b/\sqrt t)\bigr),
\]
and more precisely $\P(M_t \ge b, W_t \le a) = \P(W_t \ge 2b - a)$ for $a \le b$.
\end{theorem}

\begin{proof}
On $\{\tau_b \le t\}$ define the reflected path $\tilde W_s = W_s$ for $s \le \tau_b$ and $\tilde
W_s = 2b - W_s$ afterwards (\cref{fig:qm:brownian-motion:reflect}). By
\cref{thm:qm:brownian-motion:strongmarkov}, $W_{\tau_b + u} - b$ is a \omterm{def:qm:brownian-motion:bm}{Brownian motion}
independent of $\mathcal F_{\tau_b}$, and so is its negative; hence $\tilde W$ is a \omterm{def:qm:brownian-motion:bm}{Brownian
motion}. The event $\{M_t \ge b, W_t \le a\}$ for $W$ is the event $\{\tilde W_t \ge 2b - a\}$ for
$\tilde W$ (the paths reach $b$ at the same time), which has the probability of $\{W_t \ge 2b -
a\}$. With $a = b$: $\P(M_t \ge b) = \P(M_t \ge b, W_t \le b) + \P(W_t > b) = 2\P(W_t \ge b)$.
\end{proof}

\begin{omfigure}
\begin{tikzpicture}
\begin{axis}[width=11cm,height=5cm,
  xlabel={$t$},ylabel={$W_t$},
  tick label style={font=\small},label style={font=\small},
  xmin=0,xmax=1,ymin=-0.6,ymax=1.3,grid=major,grid style={gray!25},
  legend columns=2,legend style={font=\footnotesize,at={(0.5,-0.3)},anchor=north,draw=none,fill=none,/tikz/every even column/.append style={column sep=8pt}},
  table/col sep=comma]
\addplot[omThm,thick,dashed] table[x=t,y=refl]{figdata/methods/02-brownian-motion/reflect.csv};
\addplot[omDef,thick] table[x=t,y=w]{figdata/methods/02-brownian-motion/reflect.csv};
\addplot[black,thin,forget plot] coordinates {(0,0.6) (1,0.6)};
\addplot[gray,thin,forget plot] coordinates {(0.618,-0.6) (0.618,1.3)};
\node[font=\footnotesize,anchor=south west,fill=white,inner sep=1pt] at (axis cs:0.01,0.63) {level $b = 0.6$};
\node[font=\footnotesize,anchor=north west,fill=white,inner sep=1pt] at (axis cs:0.64,-0.33) {$\tau_b = 0.618$};
\legend{reflected after $\tau_b$,path}
\end{axis}
\end{tikzpicture}

\omcaption{The reflection principle. After the first passage at $b$, the path (solid) and its
mirror image in $b$ (dashed) are equally likely continuations, so every path that ends below
$b$ after touching it is paired with one that ends above. Data: a seeded path of the
tutorial.}\label{fig:qm:brownian-motion:reflect}
\end{omfigure}

Differentiating in $t$ gives the law of the \omterm{def:qm:brownian-motion:passage}{first-passage time}: $\tau_b$ has density
$b(2\pi t^3)^{-1/2}e^{-b^2/2t}$, finite almost surely but with infinite mean, since the density
decays like $t^{-3/2}$. A driftless price reaches any level eventually, and one should not
wait for it. The stop of the hook is the reflection principle applied to the log price.

\begin{definition}[Brownian motion with drift]\label{def:qm:brownian-motion:drift}
A \emph{Brownian motion with drift}\index{Brownian motion with drift} $\mu$ and volatility
$\sigma > 0$ is $X_t = \mu t + \sigma W_t$.
\end{definition}

\begin{proposition}[First passage with drift]\label{prop:qm:brownian-motion:driftpassage}
For $X_t = \mu t + \sigma W_t$ and $b < 0$,
\[
\P\bigl(\min_{s \le T}X_s \le b\bigr) = \Phi\Bigl(\frac{b - \mu T}{\sigma\sqrt T}\Bigr) +
e^{2\mu b/\sigma^2}\,\Phi\Bigl(\frac{b + \mu T}{\sigma\sqrt T}\Bigr).
\]
\end{proposition}

\admitted

The proof changes the measure to remove the drift and applies the reflection principle
(an exercise of chapter~5 derives it). For the hook,
$b = \ln 0.98 = -0.0202$ and $\sigma\sqrt T = 0.0866$: without drift the stop is touched with
probability $2\Phi(-0.233) = 81.6\%$; with the drift $-\sigma^2/2$ that makes the price itself
a \omterm{def:qm:probability-at-speed:martingale}{martingale}, 82.4\%.

\begin{proposition}[Discrete monitoring: the Broadie--Glasserman--Kou shift]\label{prop:qm:brownian-motion:bgk}
If $X$ is observed only at times $k\Delta t$, the probability that an observation falls at or
below $b < 0$ is, to first order in $\sqrt{\Delta t}$, the continuous probability for the
shifted level $b - \beta\sigma\sqrt{\Delta t}$, with $\beta = -\zeta(\tfrac12)/\sqrt{2\pi} \approx
0.5826$.
\end{proposition}

\admitted

The shift is the mean overshoot of the path beyond the level at the first observation below
it (Broadie, Glasserman and Kou, 1997). With daily closes, $\beta\sigma\sqrt{\Delta t} =
0.0110$: the 2\% stop behaves like a 3.1\% stop watched continuously, and is touched with
probability 71.9\%, against 72.0\% in a simulation of 200\,000 months of daily closes.
\Cref{fig:qm:brownian-motion:stops} shows the three curves over stop distances from 1\% to
10\%. The error of discrete monitoring shrinks only like $\sqrt{\Delta t}$: refining a daily
simulation to four, sixteen and sixty-four observations a day gives 77.1\%, 79.1\% and 80.5\%
against the continuous 81.6\%, the gap roughly halving each time the step is divided by
four. A simulation that wants the continuous answer should not refine the grid; it should use
the bridge of \cref{sec:qm:brownian-motion:bridge}.

\begin{omfigure}
\begin{tikzpicture}
\begin{axis}[width=11cm,height=5cm,
  xlabel={stop distance below entry (\%)},ylabel={P(stop touched in 21 days)},
  tick label style={font=\small},label style={font=\small},
  xmin=0.5,xmax=10.5,ymin=0,ymax=1,grid=major,grid style={gray!25},
  legend columns=2,legend style={font=\footnotesize,at={(0.5,-0.3)},anchor=north,draw=none,fill=none,/tikz/every even column/.append style={column sep=8pt}},
  table/col sep=comma]
\addplot[omDef,very thick] table[x=dist,y=cont]{figdata/methods/02-brownian-motion/stops.csv};
\addplot[omDef,only marks,mark=o,mark size=2.4pt] table[x=dist,y=bridge]{figdata/methods/02-brownian-motion/stops.csv};
\addplot[omThm,very thick,dashed] table[x=dist,y=bgk]{figdata/methods/02-brownian-motion/stops.csv};
\addplot[omThm,only marks,mark=*,mark size=1.6pt] table[x=dist,y=daily]{figdata/methods/02-brownian-motion/stops.csv};
\legend{continuous: reflection,daily closes + bridge (sim.),daily: shifted level,daily closes (sim.)}
\end{axis}
\end{tikzpicture}

\omcaption{Probability that a stop below the entry of a stock with 30\% volatility is touched
within 21 trading days. Watched continuously (solid) or only at the closes (dashed, the
Broadie--Glasserman--Kou shift); marks are simulations of 40\,000 months, the open ones with
the bridge crossing probability added between closes. Data: the chapter's tutorial,
seeded.}\label{fig:qm:brownian-motion:stops}
\end{omfigure}

\section{Quadratic variation}

\begin{definition}[Quadratic variation]\label{def:qm:brownian-motion:qv}
The \emph{quadratic variation}\index{quadratic variation} of a process $X$ on $[0, t]$ is
the limit in probability $[X]_t = \lim\sum_k (X_{t_{k+1}} - X_{t_k})^2$ over partitions $0 = t_0
< \dots < t_n = t$ whose mesh $\max_k(t_{k+1} - t_k)$ tends to zero, when it exists.
\end{definition}

\begin{theorem}[Quadratic variation of Brownian motion]\label{thm:qm:brownian-motion:qv}
$[W]_t = t$: the sums converge to $t$ in $L^2$, and almost surely along dyadic partitions.
Consequently the total variation $\sum_k|W_{t_{k+1}} - W_{t_k}|$ tends to infinity.
\end{theorem}

\begin{proof}
With $\Delta_k = W_{t_{k+1}} - W_{t_k}$, $\E[\sum\Delta_k^2] = \sum(t_{k+1} - t_k) = t$, and by
independence $\Var(\sum\Delta_k^2) = \sum 2(t_{k+1} - t_k)^2 \le 2t\cdot\text{mesh} \to 0$. For
dyadic partitions the variances are summable and Borel--Cantelli gives almost sure
convergence. If the total variation stayed bounded by $V$ along a sequence, then $\sum\Delta_k^2
\le V\max_k|\Delta_k| \to 0$ by uniform continuity, a contradiction.
\end{proof}

For a smooth function the sum of squared increments vanishes; for a Brownian path it is the
elapsed time, deterministically, whichever path occurs. For $\sigma W$ it is $\sigma^2 t$, so
one path observed ever more finely reveals its volatility exactly, while its drift stays
unobservable over any fixed period: the reason realised volatility (One Quant Book~1,
chapter~25) is estimated so much better than an expected return, and the reason the
microstructure noise of chapter~21 matters so much. \Cref{fig:qm:brownian-motion:qv}
samples one path at $4$ to $65\,536$ points.

\begin{omfigure}
\begin{tikzpicture}
\begin{groupplot}[group style={group size=2 by 1,horizontal sep=1.8cm},
  width=7.4cm,height=4.8cm,xmode=log,log basis x=2,xlabel={points in $[0,1]$},
  tick label style={font=\small},label style={font=\small},
  grid=major,grid style={gray!25},xmin=3,xmax=90000,
  table/col sep=comma]
\nextgroupplot[ylabel={$\sum(\Delta W)^2$},ymin=0.6,ymax=1.3]
\addplot[omDef,thick,mark=*,mark size=1.2pt] table[x=n,y=qv]{figdata/methods/02-brownian-motion/qv.csv};
\addplot[gray,dashed,domain=3:90000] {1};
\nextgroupplot[ylabel={$\sum|\Delta W|$},ymode=log,ymin=1,ymax=400]
\addplot[omThm,thick,mark=*,mark size=1.2pt] table[x=n,y=tv]{figdata/methods/02-brownian-motion/qv.csv};
\end{groupplot}
\end{tikzpicture}

\omcaption{One Brownian path on $[0,1]$ sampled at $2^k$ points. Left: the sum of squared
increments settles on $t = 1$ (1.006 at 65\,536 points). Right: the sum of absolute increments
grows like the square root of the number of points, without limit. Data: the chapter's
tutorial, seeded.}\label{fig:qm:brownian-motion:qv}
\end{omfigure}

\section{The Brownian bridge}\label{sec:qm:brownian-motion:bridge}

\begin{definition}[Brownian bridge]\label{def:qm:brownian-motion:bridge}
A \emph{Brownian bridge}\index{Brownian bridge} from $x$ to $y$ on $[t_0, t_1]$ is a \omterm{def:qm:brownian-motion:bm}{Brownian
motion} conditioned on $W_{t_0} = x$ and $W_{t_1} = y$; from $0$ to $0$ on $[0, T]$ it can be
written $W_t - (t/T)W_T$, a centred \omterm{def:qm:brownian-motion:bm}{Gaussian process} with covariance $s(T - t)/T$ for $s \le t$.
\end{definition}

\begin{proposition}[What happens between two observations]\label{prop:qm:brownian-motion:between}
Given $W_{t_0} = x$ and $W_{t_1} = y$, with $h = t_1 - t_0$: (i) $W_s$ for $t_0 < s < t_1$ is
normal with mean $x + (s - t_0)(y - x)/h$ and variance $(s - t_0)(t_1 - s)/h$, in particular
$\mathcal N((x+y)/2, h/4)$ at the midpoint; (ii) for $\sigma W$ and a level $b$ below both $x$
and $y$,
\[
\P\bigl(\min_{t_0 \le s \le t_1}\sigma W_s \le b \bigm| \sigma W_{t_0} = x, \sigma W_{t_1} = y\bigr) =
\exp\Bigl(-\frac{2(x - b)(y - b)}{\sigma^2 h}\Bigr).
\]
\end{proposition}

\begin{proof}
(i) is Gaussian conditioning of the vector $(W_{t_0}, W_s, W_{t_1})$. For (ii) take $\sigma = 1$,
$t_0 = 0$, $x = 0$ and level $-c$ with $c > 0$; by the reflection principle $\P(\min \le -c, W_h
\in dy) = \P(W_h \in -2c - dy)$, so the conditional probability is the ratio of the two normal
densities, $\exp(-((2c + y)^2 - y^2)/2h) = \exp(-2c(c + y)/h)$, which is the formula with $x - b
= c$ and $y - b = c + y$.
\end{proof}

\begin{method}[Two uses of the bridge in simulation]\label{met:qm:brownian-motion:bridge}
(i) \emph{Building paths coarse to fine}: draw $W_T = \sqrt T Z_0$, then the midpoint of each
interval from part (i), halving the intervals until the grid is fine enough. The path has the
right law, and the first normals fix its large-scale shape, which chapter~26 exploits with
low-discrepancy points. (ii) \emph{Monitoring a barrier between grid points}: simulate the path
at the observation dates only, and for each interval multiply the probabilities $1 - p_k$ of
\emph{not} crossing, with $p_k$ from part (ii). One minus the product is an unbiased estimate of
the continuous crossing probability, with no grid refinement.
\end{method}

Two closes each 1\% above the stop, one day apart at 30\% volatility, leave a 57\%
chance that the stop was touched in between: $\exp(-2 \times 0.01^2/(0.09/252))$. The
simulation of \cref{fig:qm:brownian-motion:stops} that adds these probabilities to daily
closes recovers the continuous answer, 81.6\%, with no finer grid.

\section{Tutorial: stops, reflections and quadratic variation}\label{tut:qm:brownian-motion:stops}

\begin{tutorial}
\textbf{Goal.} Simulate Brownian paths two ways, check the reflection principle and the
\omterm{def:qm:brownian-motion:qv}{quadratic variation}, and price the hook's stop-loss by continuous and discrete monitoring.
\textbf{End state:} \cref{fig:qm:brownian-motion:paths,fig:qm:brownian-motion:reflect,fig:qm:brownian-motion:stops,fig:qm:brownian-motion:qv}
and the probabilities 81.6\% and 72.0\%.
\begin{tutsteps}
\item \textbf{Paths}, by increments and by the bridge construction of \cref{met:qm:brownian-motion:bridge}(i).
\omcode{code/firm/mcengine/firm_mcengine.py}{86}{101}{Brownian paths built coarse to fine by the bridge.}\label{lst:qm:brownian-motion:bridge}
\item \textbf{The stop}: observe the log price at a given number of times a day, and
optionally add the bridge crossing probability between observations.
\omcode{code/methods/02-brownian-motion/python/qm_brownian.py}{37}{52}{Discrete monitoring of a stop, with and without the bridge correction.}\label{lst:qm:brownian-motion:touch}
\item \textbf{Run} \texttt{stop\_problem()}, \texttt{convergence\_table()} and
\texttt{fig\_brownian.py}; check that the C++20 and Rust engines draw the same normals as
Python from the same seed (\texttt{code/firm/mcengine/}).
\end{tutsteps}
\textbf{What to change next.} Replace the fixed stop by a trailing stop 2\% below the running
maximum and find its trigger probability; give the log price the drift $-\sigma^2/2$ and
compare the simulation with \cref{prop:qm:brownian-motion:driftpassage}.
\end{tutorial}

\section{Build: the Monte Carlo engine, stage one}\label{bld:qm:brownian-motion:mcengine}

\begin{build}
\bfield{Purpose} Every simulation of the miniature firm draws its paths here: this chapter's
Brownian paths, the stochastic differential equations of chapter~4, and the variance reduction
and low-discrepancy points of chapter~26.
\bfield{Interface} \texttt{SplitMix64(seed)}, \texttt{NormalStream(seed)} (Box--Muller, cosine
first); \texttt{brownian\_paths(n\_paths, n\_steps, T, seed, corr=None)};
\texttt{bridge\_paths(n\_paths, levels, T, seed)};
\texttt{bridge\_crossing\_probability(x0, x1, barrier, var)}; in C++20 and Rust,
\texttt{brownian\_path}, \texttt{bridge\_path} and \texttt{bridge\_crossing\_probability} over a
\texttt{NormalStream}.
\bfield{Rules} The reference stream (SplitMix64, Box--Muller) is identical in the three
languages; bulk Python paths use NumPy's PCG64; correlated paths through the Cholesky factor
of the correlation matrix; seeds are explicit arguments, never global state.
\bfield{Acceptance tests} \texttt{code/firm/mcengine/\{tests,cpp,rust\}}: the SplitMix64 test
vector; the first normals and a four-step path equal in the three languages to $10^{-14}$;
$\Var(W_T) = T$ and $\Cov(W_s, W_t) = \min(s,t)$ for both constructions; the crossing formula
against a fine simulation.
\bfield{Stretch} Antithetic pairs; a counter-based generator whose streams can be split across
threads (chapter~26).
\end{build}

\begin{omsources}
\item L.~Bachelier, ``Théorie de la spéculation'', \emph{Annales scientifiques de l'École
Normale Supérieure} 17, 1900.
\item N.~Wiener, ``Differential-space'', \emph{Journal of Mathematics and Physics} 2, 1923.
\item M.~D.~Donsker, ``An invariance principle for certain probability limit theorems'',
\emph{Memoirs of the AMS} 6, 1951.
\item M.~Broadie, P.~Glasserman and S.~Kou, ``A continuity correction for discrete barrier
options'', \emph{Mathematical Finance} 7, 1997.
\item I.~Karatzas and S.~E.~Shreve, \emph{Brownian Motion and Stochastic Calculus},
Springer, 2nd ed., 1991.
\end{omsources}

\section{Exercises}

\begin{exercise}[$\star$]\label{exo:qm:brownian-motion:1}
Compute $\Cov(W_2, W_5)$, $\Var(W_5 - W_2)$ and the correlation of $W_2$ and $W_5$.
\end{exercise}

\begin{exercise}[$\star$]\label{exo:qm:brownian-motion:2}
What is the probability that a standard \omterm{def:qm:brownian-motion:bm}{Brownian motion} reaches 1 before time 1?
\end{exercise}

\begin{exercise}[$\star$]\label{exo:qm:brownian-motion:3}
A stock has an annual volatility of 20\%. What is the standard deviation of its log return
over one day and over one week of five trading days?
\end{exercise}

\begin{exercise}[$\star\star$]\label{exo:qm:brownian-motion:4}
A stop sits 5\% below the entry of a stock with 30\% volatility. What is the probability it is
touched within three months (63 trading days) if watched continuously, and if watched at the
closes only?
\end{exercise}

\begin{exercise}[$\star\star$]\label{exo:qm:brownian-motion:5}
Show that $W_t^2 - t$ is a \omterm{def:qm:probability-at-speed:martingale}{martingale}, and deduce the expected time for $W$ to leave the
interval $(-a, b)$, $a, b > 0$.
\end{exercise}

\begin{exercise}[$\star\star$]\label{exo:qm:brownian-motion:6}
For a \omterm{def:qm:brownian-motion:bridge}{Brownian bridge} from 0 to 0 on $[0, 1]$, give the law of its value at $t = \tfrac12$ and the
probability that it reaches $0.5$.
\end{exercise}

\begin{exercise}[$\star\star\star$]\label{exo:qm:brownian-motion:7}
\emph{Coding.} The sum of squared increments of $W$ on $[0,1]$ at $n$ points has standard
deviation $\sqrt{2/n}$. Check it by simulating 2\,000 paths at $n = 4\,096$, and say whether the
fluctuations of \cref{fig:qm:brownian-motion:qv} are consistent with it.
\end{exercise}

\begin{exercise}[$\star\star\star$]\label{exo:qm:brownian-motion:8}
\emph{Find the flaw.} ``Our Monte Carlo for a contract that knocks out if the stock ever
trades below a barrier simulates daily closes and checks them against the barrier; the price
has converged to four digits.'' Correct it.
\end{exercise}

\section{Problem: The Stop-Loss}

\begin{problem}[{Weekend problem --- how often a 2\% stop is touched in a month}]\label{pb:qm:brownian-motion:1}
A stock with an annual volatility of 30\% is bought with a stop 2\% below the entry. Its log
price is modelled as a \omterm{def:qm:brownian-motion:bm}{Brownian motion} with volatility $\sigma = 0.30$ and no drift, over 21
trading days of a 252-day year.

\medskip\noindent\textbf{Part I --- The continuous stop.}
\begin{enumerate}
\item What is the stop level $b$ in log price, and the standard deviation $\sigma\sqrt T$ over the month?
\item What is the probability that the stock ends the month below the stop?
\item What is the probability that it touches the stop during the month, watched continuously?
\item Why is the second number exactly twice the first?
\item With the drift $-\sigma^2/2$ that makes the price a \omterm{def:qm:probability-at-speed:martingale}{martingale}, what does \cref{prop:qm:brownian-motion:driftpassage} give?
\end{enumerate}

\medskip\noindent\textbf{Part II --- Watched at the close.}
\begin{enumerate}[resume]
\item What does a simulation of 200\,000 months of daily closes give?
\item What is the Broadie--Glasserman--Kou shift with daily closes, and the equivalent continuous stop?
\item What probability does the shifted level give?
\item Watched every 15 minutes (26 times a day), what do the simulation and the shift give?
\item Why does refining the grid by four only halve the gap to the continuous answer?
\end{enumerate}

\medskip\noindent\textbf{Part III --- The bridge.}
\begin{enumerate}[resume]
\item Two consecutive closes are each 1\% above the stop. What is the probability the stop was touched between them?
\item What does the daily simulation give once each day's crossing probability is added?
\item Why is this estimate unbiased for the continuous probability?
\item Which is cheaper for a target accuracy: 64 observations a day, or daily observations with the bridge?
\end{enumerate}

\medskip\noindent\textbf{Part IV --- Judgement.}
\begin{enumerate}[resume]
\item At 15\% volatility, what is the continuous probability?
\item Over three months, what are the continuous and daily-close probabilities?
\item The stop is a resting order on the exchange; which number applies, and what does the model miss?
\item The desk sells a note that knocks out on a daily close below the level. Which number should price it?
\item State the \emph{named result}: the probability that the 2\% stop is touched in a month, continuously and at the closes.
\item In one sentence: why is a barrier watched continuously so much more likely to be touched than one watched daily?
\end{enumerate}
\end{problem}

\section{Interview questions}

\begin{interviewq}[$\star$ \iqroles{trader, researcher}]\label{iq:qm:brownian-motion:1}
A driftless \omterm{def:qm:brownian-motion:bm}{Brownian motion} starts at 0. What is the probability it hits $-1$ before $+2$?
\end{interviewq}

\begin{interviewq}[$\star$ \iqroles{researcher}]\label{iq:qm:brownian-motion:2}
Why is the \omterm{def:qm:brownian-motion:qv}{quadratic variation} of a Brownian path $t$, and why does it matter for trading?
\end{interviewq}

\begin{interviewq}[$\star\star$ \iqroles{researcher, trader}]\label{iq:qm:brownian-motion:3}
What is the expected maximum of a standard \omterm{def:qm:brownian-motion:bm}{Brownian motion} over $[0,1]$?
\end{interviewq}

\begin{interviewq}[$\star\star$ \iqroles{developer, researcher}]\label{iq:qm:brownian-motion:4}
How do you price a continuously monitored barrier by Monte Carlo without taking tiny time steps?
\end{interviewq}

\begin{interviewq}[$\star\star$ \iqroles{risk, trader}]\label{iq:qm:brownian-motion:5}
A trader says a 2\% stop on a 30\%-volatility stock is triggered about 40\% of the time in a
month. What did they compute, and what is the right number?
\end{interviewq}

\begin{interviewq}[$\star\star\star$ \iqroles{researcher}]\label{iq:qm:brownian-motion:6}
Compute $\P(W_1 > 0, W_2 > 0)$.
\end{interviewq}
